\section{The fixed-parameter algorithm}\label{sec:fpt}

The decision procedure underlying Theorem~\ref{thm:fpt} is
Algorithm~\ref{alg:fpt}. It quantifies over the $\kappa_{\rm obs}$ free-row
coordinates of \eqref{eq:kappa} and over the certificate of
Lemma~\ref{lem:window}; the only non-polynomial factor is the number
$\kappa_{\rm obs}$ of those coordinates.

\begin{algorithm}[!t]
\caption{Deciding $\mathrm{DEC}(q)$ in time $f(\kappa_{\rm obs})\operatorname{poly}(\mathcal L)$}
\label{alg:fpt}
\begin{algorithmic}[1]
\Require $A,C,L_0,q$ as above; $r=q-r_0$.
\Ensure \textsc{feasible} or \textsc{infeasible}.
\If{$r=0$}
  \State \Return the exact test of Proposition~\ref{prop:zero} on $L=L_0$
\EndIf
\State $N\gets$ basis of $\ker[C;CA]$;\quad $d_0\gets\operatorname{rank}[N\ AN]$;\quad
$\kappa_{\rm obs}\gets rd_0$
\State choose a rational basis $T$ of $\mathcal T=\operatorname{row}[N\ AN]$
\State introduce $t\in\mathbb R^{\kappa_{\rm obs}}$, the coordinates of the $r$ free rows in $T$
\State $H(t)\gets[L_0;R(t)]N$ and $H_A(t)\gets[L_0;R(t)]AN$
\For{$\omega=0$ \textbf{to} $\min(q,k)$}
  \State form $\Psi_\omega(t,x)$ with
  \Statex \quad (i) $e_\omega(G_H)>0$ and $e_{\omega+1}(G_S)=0$, where
           $G_H=H^{\mathsf T}H$ and $G_S=H^{\mathsf T}H+H_A^{\mathsf T}H_A$;
  \Statex \quad (ii) $Z^{\mathsf T}Zx=e_0$, with $Z^{\mathsf T}Z$ computed from the kernel
           $K(\lambda,\mu)$ by Cauchy--Binet (Lemma~\ref{lem:window}).
  \If{$\Psi_\omega$ is real-feasible}
    \State \Return \textsc{feasible}
  \EndIf
\EndFor
\State \Return \textsc{infeasible}
\end{algorithmic}
\end{algorithm}

\begin{thm}\label{thm:fpt}
For rational input and every $r\ge0$, $\mathrm{DEC}$ is decidable in
$f(\kappa_{\rm obs})\operatorname{poly}(\mathcal L)$ bit operations, where
$\kappa_{\rm obs}=r\operatorname{rank}[N\ AN]=r d_0$ is the effective dimension
of \eqref{eq:kappa} and $r$ is the number of free observer rows. For $r=0$ the
bound sharpens to $\operatorname{poly}(\mathcal L)$.
\end{thm}

\begin{proof}
The two cases are separated, because the free-row parametrization degenerates
at $r=0$.

\emph{Case $r=0$.} There are no free rows, hence $L=L_0$ is fixed and
$\kappa_{\rm obs}=0$, and Proposition~\ref{prop:zero} decides \eqref{eq:5a} and
\eqref{eq:5b} for this single $L$ in $\operatorname{poly}(n,\mathcal L)$ bit
operations.

\emph{Case $r\ge1$.} The $r$ free rows are parametrized by the
$\kappa_{\rm obs}$ real coordinates $t$ of \eqref{eq:kappa}. For each
$\omega=0,\dots,\min(q,k)$, form
\begin{equation*}
H(t)=[L_0;R(t)]N,\qquad H_A(t)=[L_0;R(t)]AN,
\end{equation*}
and impose
\begin{equation*}
\begin{gathered}
G_H=H^{\mathsf T}H,\qquad G_S=H^{\mathsf T}H+H_A^{\mathsf T}H_A,\\
e_\omega(G_H)>0,\qquad e_{\omega+1}(G_S)=0,
\end{gathered}
\end{equation*}
together with the $2\omega$-variable certificate
\begin{equation*}
Z^{\mathsf T}Zx=e_0,\qquad x\in\mathbb R^{2\omega},
\end{equation*}
of Lemma~\ref{lem:window}, where $Z$ is built from the maximal minors of
$P(\lambda)=\lambda H(t)-H_A(t)$ and $Z^{\mathsf T}Z$ is computed from
$K(\lambda,\mu)$ by Cauchy--Binet. At $\omega=0$ the branch is the one of
Section~\ref{sec:deg}: there $H=0$, \eqref{eq:5b} reads $H_A=0$, and
Lemma~\ref{lem:window} is not applied.

If a feasible $L=[L_0;R]$ exists, let $t$ be the coordinates of its free rows;
then $H=LN$ and $H_A=LAN$, the Gram conditions encode \eqref{eq:5a} through the
positivity of elementary symmetric functions, and the certificate encodes
\eqref{eq:5b} by Lemma~\ref{lem:window}.

Conversely, if the formula is true for some $t$, define $L$ by taking $L_0$ and,
for each free row pair in $\mathcal T$, a preimage under $[N\ AN]$; then $H=LN$
and $H_A=LAN$. The two Gram conditions give $\operatorname{rank}H=\omega$ and
$\operatorname{row}H_A\subseteq\operatorname{row}H$, hence \eqref{eq:5a} by
Lemma~\ref{lem:compression}(i), and the certificate gives \eqref{eq:5b} by
Lemma~\ref{lem:window}. The formula is therefore equivalent to feasibility.

The quantified variables and the data of the real-feasibility problem are
\begin{equation}\label{eq:nvar}
\begin{aligned}
N&=\kappa_{\rm obs}+2\omega\le3\kappa_{\rm obs}=f(\kappa_{\rm obs}),\\
s&=2\omega+O(1)=f(\kappa_{\rm obs}),\qquad
D=\kappa_{\rm obs}+1=f(\kappa_{\rm obs}),\\
B&=f(\kappa_{\rm obs})\operatorname{poly}(\mathcal L),
\end{aligned}
\end{equation}
where the discrete choice of $\omega$ is a finite disjunction, not an
existential real variable, and $B$ is the coefficient bit length of
Lemma~\ref{lem:coeff}. Substituting \eqref{eq:nvar} into the bound
\eqref{eq:rengar} of Section~\ref{sec:bitsize} gives
\begin{equation*}
\begin{gathered}
T\le B(\log B)(\log\log B)(sD)^{O(N)}\\
=\operatorname{poly}(\mathcal L)\cdot f(\kappa_{\rm obs})^{O(\kappa_{\rm obs})}\\
=f(\kappa_{\rm obs})\operatorname{poly}(\mathcal L),
\end{gathered}
\end{equation*}
bit operations, which is the stated bound.
\end{proof}

\begin{cor}\label{cor:arith}
The same algorithm uses $f(\kappa_{\rm obs})\operatorname{poly}(\mathcal L)$
arithmetic operations.
\end{cor}

\subsection{The bit model and the coefficient bound}\label{sec:bitsize}

We use the standard bit model: the input has rational entries of total bit
length $\mathcal L$, and $n,p,q\le n\le\mathcal L$. Renegar's Part II bit bound
\cite{ren92} has the form
\begin{equation}\label{eq:rengar}
B(\log B)(\log\log B)(sD)^{O(N)},
\end{equation}
for integer-coefficient input, where $N$ is the number of real variables, $s$
the number of polynomials, $D$ the maximum degree, and $B$ the maximum
coefficient bit length. The exponential factor depends only on $N$; the
coefficient-length contribution is polynomial in $\log B$.

\begin{lem}[Coefficient sizes]\label{lem:coeff}
Let the input entries have bit length at most $\mathcal L$, and let
$n\le\mathcal L$. Then:
\emph{(i)} the kernel basis $N$, the matrix $D$, and the entries of
$H,H_A,G_H,G_S$ have bit length $\operatorname{poly}(n,\mathcal L)$;
\emph{(ii)} every maximal minor $f_i$ of $P(\lambda)=\lambda H-H_A$, and every
coefficient of $f_i$, has bit length $\operatorname{poly}(n,\mathcal L)$;
\emph{(iii)} the entries of $Z^{\mathsf T}Z$ and of $e_j(G_H),e_j(G_S)$ for
$j\le\kappa_{\rm obs}+1$ have bit length
$f(\kappa_{\rm obs})\operatorname{poly}(\mathcal L)$.
\end{lem}

\begin{proof}
Let $m=n^2+pn+r_0n$ be the number of rational input entries. Multiplying by the
product of their denominators turns the input into an integer matrix
$A_{\rm in}$ all of whose entries are bounded by $B_0$ in absolute value, where
\begin{equation*}
m\le3n^2=O(n^2),\qquad B_0=2^{O(m\mathcal L)}=2^{O(n^2\mathcal L)} .
\end{equation*}
Hadamard's inequality gives, for a $t\times t$ minor $M_{I,J}$ of $A_{\rm in}$
with $t\le n$,
\begin{equation*}
\begin{aligned}
|\det M_{I,J}|&\le(\sqrt t\,B_0)^t,\\
\operatorname{bits}\det M_{I,J}
&=O\bigl(t(\log B_0+\log t)\bigr)=O(n^3\mathcal L),
\end{aligned}
\end{equation*}
which is $\operatorname{poly}(n,\mathcal L)$ because $n\le\mathcal L$.

\emph{(i)} The kernel basis $N$ of $[C;CA]$, the row-space basis of $[N\ AN]$,
and hence the matrix $D$ are obtained by Cramer's rule or by reduced row echelon
form, so every entry is a ratio of such minors and the common denominator is a
product of at most $n$ pivot minors:
\begin{equation*}
\operatorname{bits}N_{ij},\ \operatorname{bits}D_{ij}
\le 2\cdot O(n^3\mathcal L)+O(\log n)=\operatorname{poly}(n,\mathcal L).
\end{equation*}
Each entry of $H=LN$ is a sum of $n$ products of an input entry with an entry of
$N$, and each entry of $G_H$ or $G_S$ is a sum of at most $q$ such products, so
passing from $N$ to $H$ and then to $G_H,G_S$ adds $O(\log n)$ and then
$O(\log q)$ bits:
\begin{equation*}
\begin{aligned}
\operatorname{bits}H_{ij}&\le\mathcal L+\operatorname{poly}(n,\mathcal L)+O(\log n),\\
\operatorname{bits}(G_H)_{ij},\ \operatorname{bits}(G_S)_{ij}
&\le2\operatorname{poly}(n,\mathcal L)+O(\log q).
\end{aligned}
\end{equation*}

\emph{(ii)} A maximal minor $f_i$ of $P(\lambda)=\lambda H-H_A$ is a determinant
of an $\omega\times\omega$ matrix whose entries are linear in those of $H,H_A$;
Hadamard applied to its coefficient matrix bounds each of its $\omega+1$
coefficients:
\begin{equation*}
\deg f_i\le\omega,\qquad
\operatorname{bits}c_{i,a}=\operatorname{poly}(n,\mathcal L).
\end{equation*}

\emph{(iii)} $Z^{\mathsf T}Z$ is a sum of at most
$\binom{k}{\omega}\le2^{\kappa_{\rm obs}}$ terms, each a product of two
coefficients of the $f_i$, and $e_j$ is a sum of at most
$k^{j}\le\kappa_{\rm obs}^{\kappa_{\rm obs}}$ principal minors. A sum of $T$
terms of bit length $b$ has bit length $b+\log T$, hence
\begin{equation*}
\begin{aligned}
\operatorname{bits}(Z^{\mathsf T}Z)_{ij}
&\le2\operatorname{poly}(n,\mathcal L)+\kappa_{\rm obs},\\
\operatorname{bits}e_j(G_H),\ \operatorname{bits}e_j(G_S)
&\le\operatorname{poly}(n,\mathcal L)+\kappa_{\rm obs}\log\kappa_{\rm obs},
\end{aligned}
\end{equation*}
which is $f(\kappa_{\rm obs})\operatorname{poly}(\mathcal L)$.
\end{proof}

It remains to count the cost of Algorithm~\ref{alg:fpt} at fixed
$\kappa_{\rm obs}$. Assume $r\ge1$, so that
\begin{equation*}
k\le d_0\le\kappa_{\rm obs},\qquad j\le\omega+1\le\kappa_{\rm obs}+1,
\end{equation*}
the two inequalities being elementary rank-nullity bounds. The Gram matrices
$G_H$ and $G_S$ are $k\times k$, hence of size at most
$\kappa_{\rm obs}$, and their elementary symmetric polynomials are evaluated by
determinant expansion in $k^{O(j)}=f(\kappa_{\rm obs})$ operations. The matrix
$Z^{\mathsf T}Z$ is computed from the Cauchy--Binet kernel $K(\lambda,\mu)$: the
sum has at most $\binom{k}{\omega}\le2^{\kappa_{\rm obs}}$ terms, each an
$\omega\times\omega$ determinant, and all coefficients are rational of bit length
$f(\kappa_{\rm obs})\operatorname{poly}(\mathcal L)$ by Lemma~\ref{lem:coeff}.
The preprocessing therefore costs
$f(\kappa_{\rm obs})\operatorname{poly}(\mathcal L)$ bit operations. The
real-feasibility problem itself has
\begin{equation*}
\begin{gathered}
N=\kappa_{\rm obs}+2\omega=f(\kappa_{\rm obs})\ \text{ variables},\\
s\le2\omega+O(1)=f(\kappa_{\rm obs})\ \text{ equations},
\end{gathered}
\end{equation*}
degree $D\le\kappa_{\rm obs}+1$, and coefficient bit length
$B=f(\kappa_{\rm obs})\operatorname{poly}(\mathcal L)$, all as in
\eqref{eq:nvar}; by \eqref{eq:rengar},
\begin{equation*}
B(\log B)(\log\log B)(sD)^{O(N)}
=f(\kappa_{\rm obs})\operatorname{poly}(\mathcal L)
\end{equation*}
bit operations. The cases $k=0$ and $\omega=0$ are direct, and $r=0$ is decided
separately by Proposition~\ref{prop:zero}; all degenerate cases are collected in
Section~\ref{sec:deg}. For fixed $q$ the sharper parameter is $q$ itself, and the
bound $\mathcal L^{O(q^2)}$ of Section~\ref{sec:fixedq} applies instead.

\subsection{The zero-free-row case: an exact polynomial-time pencil test}\label{sec:zero}

When $r=0$ there are no free rows: $L=L_0$ is fixed and $\kappa_{\rm obs}=0$.
Only two conditions remain, and both are decidable without a
parameter-dependent search. \eqref{eq:5a} is a constant-rank test on
$H=L_0N$, $H_A=L_0AN$ (Lemma~\ref{lem:compression}(i)); \eqref{eq:5b} is the
constancy of the rank of the degree-one pencil $\lambda X-G$. We first isolate
the pencil statement.

\begin{lem}\label{lem:zero}
Let $X,G\in\mathbb Q^{q\times\omega}$ with $q\ge\omega\ge1$, and put
$Q(\lambda)=\lambda X-G$. The following are equivalent:
\emph{(i)} $\operatorname{rank}Q(\lambda)=\omega$ for every
$\lambda\in\mathbb C$;
\emph{(ii)} the maximal $\omega\times\omega$ minors of $Q$ have no common
complex root;
\emph{(iii)} the maximal minors do not all vanish and their gcd in
$\mathbb Q[\lambda]$ is a nonzero constant;
\emph{(iv)} there is $S(\lambda)\in\mathbb Q[\lambda]^{\omega\times q}$ with
$S(\lambda)Q(\lambda)=I_\omega$ and $\deg S\le2\omega-2$;
\emph{(v)} there exist matrices
$S_0,\dots,S_{2\omega-2}\in\mathbb Q^{\omega\times q}$ satisfying the rational
linear system
\begin{equation}
\tag{R0}\label{eq:R0}
\begin{aligned}
-S_0G&=I_\omega,\\
S_{j-1}X-S_jG&=0\quad(1\le j\le2\omega-2),\\
S_{2\omega-2}X&=0 .
\end{aligned}
\end{equation}
\end{lem}

\begin{proof}
\emph{(i)$\Leftrightarrow$(ii).} Since $Q$ is $q\times\omega$,
\begin{equation*}
\operatorname{rank}Q(\lambda)<\omega
\iff f_I(\lambda)=0\ \text{ for every } I\subseteq[q],\ |I|=\omega,
\end{equation*}
so a rank drop occurs exactly when all maximal minors vanish at the same point.

\emph{(ii)$\Leftrightarrow$(iii).} The minors have rational coefficients and
$\mathbb C$ is algebraically closed. A nonconstant gcd $g$ has a root
$\alpha\in\mathbb C$, and $g\mid f_I$ gives $f_I(\alpha)=0$ for all $I$;
conversely, if $f_I(\alpha)=0$ for all $I$, the minimal polynomial
$m_\alpha\in\mathbb Q[\lambda]$ of $\alpha$ divides every $f_I$ and hence
divides the gcd. Thus
\begin{equation*}
\begin{aligned}
\gcd_I f_I\ \text{nonconstant}
&\Longrightarrow\exists\alpha\in\mathbb C:\ f_I(\alpha)=0\ \ \forall I\\
&\Longrightarrow m_\alpha\mid\gcd_I f_I ,
\end{aligned}
\end{equation*}
with each implication reversible, and the clause that the minors do not all
vanish excludes the zero gcd.

\emph{(iii)$\Rightarrow$(iv).} This is the adjugate/B\'ezout construction of
Appendix~\ref{app:encoding}. For each $\omega$-element row set
$I\subseteq[q]$, let $E_I\in\mathbb Q^{\omega\times q}$ be the row selection,
$Q_I=E_IQ$ and $f_I=\det Q_I$. Then
\begin{equation*}
\operatorname{adj}(Q_I)E_IQ=f_II_\omega,\qquad
\sum_Ig_If_I=1\ \text{ with }\ \deg g_I<\omega
\end{equation*}
by Lemma~\ref{lem:bezout}, so that
\begin{equation*}
S=\sum_Ig_I\operatorname{adj}(Q_I)E_I
\quad\text{satisfies}\quad
S(\lambda)Q(\lambda)=I_\omega ,
\end{equation*}
and $\deg\operatorname{adj}(Q_I)\le\omega-1$ gives
$\deg S\le(\omega-1)+(\omega-1)=2\omega-2$.

\emph{(iv)$\Rightarrow$(v).} Write $S=\sum_{j=0}^{2\omega-2}S_j\lambda^j$. Then
$S(\lambda)(\lambda X-G)=I_\omega$, and comparing the coefficients of
$\lambda^0$, $\lambda^k$ ($1\le k\le2\omega-2$) and $\lambda^{2\omega-1}$ gives
\begin{equation*}
-S_0G=I_\omega,\qquad S_{k-1}X-S_kG=0,\qquad S_{2\omega-2}X=0,
\end{equation*}
which is \eqref{eq:R0}.

\emph{(v)$\Rightarrow$(i).} Evaluating the polynomial identity
$S(\lambda)Q(\lambda)=I_\omega$ at any $\lambda_0\in\mathbb C$ exhibits a left
inverse of $Q(\lambda_0)$:
\begin{equation*}
S(\lambda_0)Q(\lambda_0)=I_\omega
\Longrightarrow\operatorname{rank}Q(\lambda_0)\ge\omega,
\end{equation*}
while $\operatorname{rank}Q(\lambda_0)\le\omega$ because $Q(\lambda_0)$ has
$\omega$ columns; hence $\operatorname{rank}Q(\lambda_0)=\omega$.
\end{proof}

The size of \eqref{eq:R0} is fixed by the degree bound $\deg S\le2\omega-2$.
With $d=2\omega-2$ the unknowns are the entries of $d+1=2\omega-1$ matrices of
size $\omega\times q$, and each of the $d+2=2\omega$ matrix equations contributes
$\omega^2$ scalar equations, so
\begin{equation*}
N_{\rm var}=(2\omega-1)q\omega=O(q\omega^2),\qquad
N_{\rm eq}=(2\omega)\omega^2=2\omega^3 .
\end{equation*}
All coefficients of \eqref{eq:R0} are entries of $X$, $G$, or $0,\pm1$, so
clearing the denominators multiplies the system by an integer of bit length
$O(\mathcal L_0)$, where $\mathcal L_0$ is the input length of $(X,G)$. Exact
fraction-free elimination on the resulting integer system of size
$M=N_{\rm var}+N_{\rm eq}=O(\omega^3+q\omega^2)$ then satisfies
\begin{equation*}
\begin{aligned}
\text{entry bits}&=O(\mathcal L_0),\\
\text{minor bits}&=O\bigl(M(\mathcal L_0+\log M)\bigr)\ \text{(Hadamard)},\\
\text{arithmetic operations}&=O(M^3),\\
\text{bit operations}&=O(M^3)\cdot O\bigl(M(\mathcal L_0+\log M)\bigr)^2\\
&=\operatorname{poly}(M,\mathcal L_0).
\end{aligned}
\end{equation*}
Hence consistency of \eqref{eq:R0} is decided in
$\operatorname{poly}(q,\omega,\mathcal L_0)$ bit operations.

\begin{remark}[Construction and algorithm]\label{rem:construction}
The adjugate construction in (iii)$\Rightarrow$(iv) of Lemma~\ref{lem:zero}
enumerates the $\binom q\omega$ maximal minors, which is exponential in $q$; it
is used only to prove the existence of a certificate of degree at most
$2\omega-2$. It applies to the rectangular pencil $Q$ because each $Q_I$ is
square, so that $\deg S\le2\omega-2$ holds for every $q\ge\omega$ with no square
assumption on $Q$. The algorithm never forms the sum: it searches for the
coefficients of $S$ directly in \eqref{eq:R0}, whose size is polynomial in
$q,\omega$. This is the step that turns a certificate-size statement into a
decision algorithm.
\end{remark}

\begin{remark}[Higher degree]\label{rem:degree}
The same argument applies to pencils of degree $d$: if the entries of $Q$ have
degree at most $d$, then every maximal minor has degree at most $d\omega$,
B\'ezout gives multipliers with $\deg g_I<d\omega$, and the adjugate contributes
at most $d(\omega-1)$,
\begin{equation*}
\deg S\le(d\omega-1)+d(\omega-1)=2d\omega-d-1,
\end{equation*}
which is $2\omega-2$ for $d=1$. Degree bounds for polynomial-matrix left
inverses are studied classically \cite{fab90}; the bound above is the elementary
adjugate--B\'ezout one, and only the case $d=1$ is used in this paper. The
general bound may be useful if the observer model is extended. For $d=0$ the
pencil is constant and the statement is the rank test of
Lemma~\ref{lem:compression}(iii).
\end{remark}

\begin{prop}\label{prop:zero}
Let $r=0$, so $L=L_0$. Feasibility of $\eqref{eq:5a}\wedge\eqref{eq:5b}$ is
decidable in $\operatorname{poly}(n,\mathcal L)$ bit operations.
\end{prop}

\begin{proof}
Compute a basis $N$ of $\ker[C;CA]$ and set $H=L_0N$, $H_A=L_0AN$.

\emph{Step 1.} Check $\operatorname{rank}[H_A;H]=\operatorname{rank}H$; if it
fails, reject, since this is \eqref{eq:5a}
(Lemma~\ref{lem:compression}(i)).

\emph{Step 2.} Put $\omega=\operatorname{rank}H$. If $\omega=0$, then $H=0$ and
\eqref{eq:5b} reads $\operatorname{rank}(\lambda H-H_A)=0$ for all $\lambda$,
i.e.\ $H_A=0$; accept exactly in that case.

\emph{Step 3.} Otherwise take a rank factorization $H=XB$ with
$X\in\mathbb Q^{q\times\omega}$ of full column rank and
$B\in\mathbb Q^{\omega\times k}$ of full row rank. Step 1 gives
$\operatorname{row}H_A\subseteq\operatorname{row}H$
(Lemma~\ref{lem:compression}(i)), hence $H_A=FH$ for some $F$; set $G=FX$, so
$H_A=GB$.

\emph{Step 4.} Run the test of Lemma~\ref{lem:zero} on $(X,G)$.

\emph{Correctness.} Since $B$ is of full row rank,
\begin{equation*}
\operatorname{rank}(\lambda H-H_A)
=\operatorname{rank}\bigl((\lambda X-G)B\bigr)
=\operatorname{rank}(\lambda X-G)
\end{equation*}
for every $\lambda$, and \eqref{eq:5b} is
$\operatorname{rank}(\lambda H-H_A)=\omega$ for all $\lambda$; thus
Lemma~\ref{lem:zero} decides \eqref{eq:5b}. Step 1 is \eqref{eq:5a}, and step 2
is the exact $\omega=0$ reading of the same identity. The zero-free-row decision
is therefore the conjunction of one constant-rank test and one pencil test: if
\eqref{eq:5a} fails, the instance is infeasible and the spectral question is not
asked.

\emph{Complexity.} Step 1 is one rank computation, step 3 a rank factorization
and an exact linear solve, and step 4 costs
$\operatorname{poly}(q,\omega,\mathcal L_0)\subseteq\operatorname{poly}(n,\mathcal L)$
by Lemma~\ref{lem:zero}. Every step is a rank computation or an exact linear
solve on a matrix of at most $2\omega+1$ rows and $q$ columns whose entries have
bit length $\operatorname{poly}(n,\mathcal L)$ by Lemma~\ref{lem:coeff}, and
Gaussian elimination over $\mathbb Q$ on such a matrix costs
\begin{equation*}
\operatorname{poly}(n,\mathcal L)\ \text{ bit operations};
\end{equation*}
the test returns the coefficients $S_0,\dots,S_{2\omega-2}$ explicitly.
\end{proof}

\begin{remark}[Kronecker canonical form]\label{rem:kronecker}
The pencil $Q$ has full column rank at every finite complex value if and only if
its Kronecker form (in the sense of \cite{bee88}) contains no right singular
block and no finite regular elementary-divisor block. Left singular blocks and
\emph{infinite} regular elementary divisors are permitted: they do not produce a
rank drop at any finite $\lambda$. For instance
\begin{equation*}
Q(\lambda)=1,\qquad Q(\lambda)=I_2,\qquad
Q(\lambda)=\begin{bmatrix}1&\lambda\\0&1\end{bmatrix}
\end{equation*}
all satisfy the condition, the first with an infinite elementary divisor of size
one and the last with determinant the constant $1$. Equivalently, in the
language of Appendix~\ref{app:encoding}, the condition is that the ideal
generated by the maximal minors is the unit ideal; only the \emph{finite}
spectrum is constrained.
\end{remark}

\begin{remark}[Witness]\label{rem:witness}
If \eqref{eq:R0} is consistent, exact elimination returns
$S_0,\dots,S_{2\omega-2}$, hence an explicit polynomial left inverse
$S(\lambda)$, verifiable by multiplying the two polynomial matrices. The test is
therefore not merely a yes/no decision.
\end{remark}
